\SetPractica{Compuestos orgánicos e inorgánicos}{Diferencias fisicoquímicas entre compuestos orgánicos e inorgánicos}

\section{Introducción teórica}

Los compuestos orgánicos e inorgánicos presentan diferencias fisicoquímicas. Comenzando estas tenemos la solubilidad, que se define como la capacidad de una sustancia para disolverse en otra. La sustancia a disolver se conoce como soluto, mientras que el medio en que el solito se disuelve se llama disolvente. La solubilidad está altamente relacionada a la polaridad, pues la interacción molecular que permite la separación de las moléculas del soluto para unirse a las moléculas del disolvente se basa en la polaridad de las moléculas de los compuestos. La relación entre la polaridad de los compuestos debe ser homóloga para que la disolución ocurra, osea, un soluto polar necesita de un disolvente polar, así como un soluto no polar necesita de un disolvente no polar. Si la polaridad de los reactivos es diferente entre sí la disolución no ocurre.

Las respuestas de los compuestos a la aplicación de calor son también propiedades afectadas por la polaridad de las moléculas. Los puntos de ebullición de los compuestos tienden a aumentar cuando presentan las fuerzas intermoleculares que mencionamos previamente. Así pues, cuánto más polar es un compuesto, presenta también fuerzas intermoleculares más fuertes que necesitan de más energía para que se de la evaporación, lo que le da un valor más alto a su punto de ebullición.\citep{troChemistryStructureProperties2015}

\section{Objetivos}

\begin{itemize}
    \item Compruebe las diferencias entre los compuestos orgánicos e inorgánicos.
    \item Conozca algunas características de estos compuestos.
\end{itemize}

\section{Materiales, equipos y reactivos}

\begin{table}[H]
{\small\setstretch{1.4}
\begin{tabular}{|m{0.325\linewidth}|m{0.325\linewidth}|m{0.325\linewidth}|} \hline
    
    \thead{Equipos} & \thead{Materiales} & \thead{Reactivos} \\
    
    {Mechero Bunsen \par}
    {Soporte para mechero \par}
    {Tela de asbesto \par}
    {Fósforos o encendedor \par}
    {Termómetro \par}
    {Balanza}
    
    &
    
    {Tubos de ensayo \par}
    {Gradilla \par}
    {Pinzas \par}
    {Agitador \par}
    {Vasos de precipitado \par}
    {Pipetas \par}
    {Cápsulas de porcelana \par}
    {Espátula \par}
    {Embudo}
    
    &
    
    {Agua destilada \par}
    {Cloruro de sodio (\ce{NaCl}) \par}
    {Almidón \par}
    {Azúcar \par}
    {Alcohol etílico \par}
    {Aceite}
    
    \\ \hline
\end{tabular}}\end{table}

\section{Procedimientos}

\subsection{Solubilidad}

\begin{enumerate}
    \item Colocar 2mL de agua destilada en cada tubo de ensayo y enumerarlos.
    \item Agregar una pequeña cantidad de cada reactivo.
    \item Agitar
    \item Anotar sus observaciones en una tabla.
\end{enumerate}

\subsection{Punto de ebullición} 

\begin{enumerate}
    \item Colocar en un vaso de precipitado 10mL de agua destilada y en otro 10mL de alcohol.
    \item Calentar con cuidado y anotar la temperatura de ebullición con un termómetro.
    \item Anote sus observaciones en una tabla.
\end{enumerate}

\subsection{Formación de carbono}

\begin{enumerate}
    \item En una cápsula de porcelana quema un trozo de papel periódico y anote sus observaciones.
    \item En otra cápsula calentar un azúcar y luego cloruro de sodio.
    \item Anote sus observaciones en una tabla.
\end{enumerate}

\subsection{Estabilidad térmica}

\begin{enumerate}
    \item Colocar en un tubo de ensayo 1 g de NaCl y en otro 1 g de almidón.
    \item Llevarlos a la flama del mechero y calentarlos hasta que se observe un cambio.
    \item Anote sus observaciones en una tabla.
\end{enumerate}

\bibliography{Main}